% Part: proof-theory
% Chapter: sequent-calculus
% Section: proof-examples

\documentclass[../../../include/open-logic-section]{subfiles}

\begin{document}

\olfileid[tr]{pt}{seq}{ex}

\olsection{İspat Örnekleri}

Sekantların !!{proof}s{}ını vermek için genellikle son sekanttan yukarı doğru
ilerlemek elverişlidir: içindeki bir !!a{formula} etkin !!{formula} olarak
seçilir, karşılık gelen kuralın öncülleri sekantın üstüne yazılır ve
aksiyomlara ulaşılıncaya dek işlem yinelenir. Örneğin şu sekantı ispatlamak
istediğimizi varsayalım:
\[(!C \land !D) \lif !E \Sequent !C \lif (!D \lif !E)\]
$!C \lif (!D \lif !E)$'yi ele alın: $!A \lif !B$ biçimindedir ve sekantın
sağında geçer. Bütün sekantı \Log{G1c}'nin $\RightR{\lif}$ kuralıyla
eşleştirelim:
\[\Axiom$ !A, \Gamma \fCenter \Delta, !B$
\RightLabel{\RightR{\lif}}
\UnaryInf$ \Gamma \fCenter \Delta, !A \lif !B $
\DisplayProof\]
Bu eşleştirme $\Gamma = \{(!C \land !D) \lif !E\}$, $\Delta = \emptyset$,
$!A \ident !C$ ve $!B \ident !D \lif !E$ almamız gerektiğini düşündürür.
Öyleyse bir ispatın son çıkarımı şöyle olabilir:
\[\Axiom$ !C, (!C \land !D) \lif !E \fCenter !D \lif !E$
\RightLabel{\RightR{\lif}}
\UnaryInf$(!C \land !D) \lif !E \fCenter !C \lif (!D \lif !E)$
\DisplayProof\]
$!D \lif !E$'yi etkin formül seçerek aynı düşünceyi yinelersek şunu
üretiriz:
\[\Axiom$ !D, !C, (!C \land !D) \lif !E \fCenter !E$
\RightLabel{\RightR{\lif}}
\UnaryInf$!C, (!C \land !D) \lif !E \fCenter !D \lif !E$
\RightLabel{\RightR{\lif}}
\UnaryInf$(!C \land !D) \lif !E \fCenter !C \lif (!D \lif !E)$
\DisplayProof\]
Şimdi $(!C \land !D) \lif !E$'ye odaklanıp \Log{G1c}'nin
$\LeftR{\lif}$ kuralını kullanmalıyız:
\[\Axiom$ \Gamma \fCenter \Delta, !A$
\Axiom$ !B, \Gamma \fCenter \Delta$
\RightLabel{\LeftR{\lif}}
\BinaryInf$ !A \lif !B, \Gamma \fCenter \Delta$
\DisplayProof\]
Bu bizi şuraya götürür:
\[\Axiom$!D, !C \fCenter !E, !C \land !D$
\Axiom$!D, !C, !E \fCenter !E$
\RightLabel{\LeftR{\lif}}
\BinaryInf$ !D, !C, (!C \land !D) \lif !E \fCenter !E$
\RightLabel{\RightR{\lif}}
\UnaryInf$!C, (!C \land !D) \lif !E \fCenter !D \lif !E$
\RightLabel{\RightR{\lif}}
\UnaryInf$(!C \land !D) \lif !E \fCenter !C \lif (!D \lif !E)$
\DisplayProof\]
Sol üst sekantta $!C \land !D$ !!{formula}{}ünü başlıca formül seçip
$\RightR{\land}$ kuralını kullanırsak şunu elde ederiz:
\[
\Axiom$!D, !C \fCenter !E, !C$
\Axiom$!D, !C \fCenter !E, !D$
\RightLabel{\RightR{\land}}
\BinaryInf$!D, !C \fCenter !E, !C \land !D$
\Axiom$!D, !C, !E \fCenter !E$
\RightLabel{\LeftR{\lif}}
\BinaryInf$ !D, !C, (!C \land !D) \lif !E \fCenter !E$
\RightLabel{\RightR{\lif}}
\UnaryInf$!C, (!C \land !D) \lif !E \fCenter !D \lif !E$
\RightLabel{\RightR{\lif}}
\UnaryInf$(!C \land !D) \lif !E \fCenter !C \lif (!D \lif !E)$
\DisplayProof
\]
Buraya kadar her şey yolunda. Şimdiye dek kullandığımız kurallar \Log{G1c}
ile~\Log{G3c}'de aynı olduğundan yaptıklarımızın tümü \Log{G3c}'de de
işler. En üstteki bütün sekantların solunda ve sağında aynı !!{formula}
bulunan bir !!a{proof} parçasına ulaştık. Ancak bunlar henüz tam olarak
aksiyom değildir.

\Log{G1c}'de en üstteki sekantların $!A \Sequent !A$ biçimindeki
aksiyomlardan !!{proof}s{}ını vermeliyiz. Bu, zayıflatma kurallarıyla kolayca
yapılır. Örneğin en soldaki $!D, !C \Sequent !E, !C$ üst sekantı şöyle
!!{prove}{}nabilir:
\[
\Axiom$!C \fCenter !C$
\RightLabel{\LeftR{\Weakening}}
\UnaryInf$!D, !C \fCenter !C$
\RightLabel{\RightR{\Weakening}}
\UnaryInf$!D, !C \fCenter !E, !C$
\]
Genel olarak \Log{G1c}'de şu görünüme sahip bir !!a{proof}{}ımız vardır:
\[
\Axiom$!C \fCenter !C$
\RightLabel{\LeftR{\Weakening}}
\UnaryInf$!D, !C \fCenter !C$
\RightLabel{\RightR{\Weakening}}
\UnaryInf$!D, !C \fCenter !E, !C$
\Axiom$!D \fCenter !D$
\RightLabel{\LeftR{\Weakening}}
\UnaryInf$!D, !C \fCenter !D$
\RightLabel{\RightR{\Weakening}}
\UnaryInf$!D, !C \fCenter !E, !D$
\RightLabel{\RightR{\land}}
\BinaryInf$!D, !C \fCenter !E, !C \land !D$
\Axiom$!E \fCenter !E$
\RightLabel{\LeftR{\Weakening}}
\UnaryInf$!C, !E \fCenter !E$
\RightLabel{\LeftR{\Weakening}}
\UnaryInf$!D, !C, !E \fCenter !E$
\RightLabel{\LeftR{\lif}}
\BinaryInf$ !D, !C, (!C \land !D) \lif !E \fCenter !E$
\RightLabel{\RightR{\lif}}
\UnaryInf$!C, (!C \land !D) \lif !E \fCenter !D \lif !E$
\RightLabel{\RightR{\lif}}
\UnaryInf$(!C \land !D) \lif !E \fCenter !C \lif (!D \lif !E)$
\DisplayProof
\]
Bu ispatın yapısı (özellikle yüksekliği), $!C$, $!D$ ve~$!E$
!!{formula}s{}inin ne olduğundan bağımsızdır. Yukarıdaki !!{proof} içinde bu
simgeleri gelişigüzel !!{formula}s ile değiştirebiliriz ve sonuç yine
\Log{G1c}'de bir !!a{proof} olur. \Log{G1c}'deki ispat \emph{şematiktir}.

\Log{G3c}'de aksiyomlar, $!A$ atomik olmak koşuluyla $!A, \Gamma \Sequent
\Delta, !A$ biçimindeki sekantlardır. Dolayısıyla $!D, !C \Sequent !E, !C$
bir \Log{G3c} aksiyomu gibi görünse de ($\Gamma = \{!D\}$; $\Delta =
\{!E\}$) ancak $!C$ gerçekten atomikse bir aksiyomdur. İspat parçamız
yalnız $!C$, $!D$ ve $!E$ atomik !!{formula}s olduğunda \Log{G3c}'de tam
bir ispattır.

Tam anlamıyla
\[\Log{G1c} \Proves (!C \land !D) \lif !E \fCenter !C \lif (!D \lif
!E)\]
olduğunu göstermemiş olsak da uygulamada yapmamız gereken (ve
yapabileceğimiz) bundan ibarettir. Çünkü $!A$ atomik olmasa bile $!A,
\Gamma \Sequent \Delta, !A$ biçimindeki her sekant \Log{G3c}'de
ispatlanabilir. Bunu $!A$'nın derinliği~$\depth{!A}$ üzerinde tümevarımla
gösterebiliriz.

\begin{defn}\ollabel{defn:depth}
Bir !!a{formula}{}ün \emph{derinliği}~$\depth{!A}$ şöyle tanımlanır:
\begin{enumerate}
  \item $!A$ atomikse $\depth{!A} = 0$.
  \item $!A \ident \lnot !B$, $!A \ident \lforall[x][!B]$ ya da
  $!A \ident \lexists[x][!B]$ ise $\depth{!A} = \depth{!B} + 1$.
  \item $!A \ident (!B \land !C)$, $!A \ident (!B \lor !C)$ ya da
  $!A \ident (!B \lif !C)$ ise $\depth{!A} =
  \max(\depth{!B},\depth{!C}) + 1$.
\end{enumerate}
\end{defn}

Her $t$ terimi için $\depth{A(x)} = \depth{A(t)}$ olduğunu ispatlamak zor
değildir. Bizim için önemli olgu şudur: her kuralda başlıca !!{formula}{}ün
derinliği, etkin !!{formula}{}ün derinliğinden her zaman büyüktür. Örneğin:
\begin{align*}
\depth{P(x)} = \depth{P(a)} & = 0,\\
\depth{\lforall[x][P(x)]} & = 1,\\
\depth{(\lforall[x][P(x)] \lif P(a))} & = \max(1,0) + 1 = 2.
\end{align*}
Bunu, aşağıdaki gibi bazı sonuçları $\depth{!A}$ üzerinde tümevarımla
ispatlamak için kullanabiliriz.

\begin{prop}\ollabel{prop:G1c-axioms}
Her~$!A$ için \Log{G1c}'de $!A \Sequent !A$'nın bütün aksiyomları atomik
olan bir !!a{proof}{}ı vardır.
\end{prop}

\begin{proof}
  $\depth{!A}$ üzerinde tümevarımla. $\depth{!A} = 0$ ise $!A$ atomiktir
  ve $!A \Sequent !A$ zaten \Log{G1c}'de istenen biçimde bir !!a{proof}{}tır.
  $\depth{!A} > 0$ ise $!A$, daha düşük derinlikteki !!{formula}s{}den
  kurulmuştur; örneğin $!A \ident \lnot !B$, $!A \ident (!B \land !C)$
  ya da $!A \ident \lforall[x][A(x)]$ olabilir. Tümevarım varsayımı,
  $\depth{!D} < \depth{!A}$ olan bütün $!D$ için $\Log{G1c} \Proves !D
  \Sequent !D$ eşitliğinin atomik aksiyomlardan elde edildiğidir.

  $!A \ident \lnot !B$ ise $\depth{!B} < \depth{!A}$ ve tümevarım
  varsayımına göre \Log{G1c}'de atomik aksiyomlardan $!B \Sequent !B$'nin
  bir !!a{proof}{}ı~$\pi_1$ vardır. Bunu şöyle $\lnot !B \Sequent \lnot !B$'nin
  bir !!a{proof}{}a genişletebiliriz:
  \[
  \AxiomC{}
  \RightLabel{$\pi_1$}
  \Deduce$!B \fCenter !B$
  \RightLabel{\LeftR{\lnot}}
  \UnaryInf$\lnot !B, !B \fCenter $
  \RightLabel{\RightR{\lnot}}
  \UnaryInf$\lnot !B \fCenter \lnot !B$
  \DisplayProof
\]

  $!A \ident (!B \land !C)$ ise $\depth{!B} < \depth{!A}$ ve
  $\depth{!C} < \depth{!A}$'dır. Tümevarım varsayımına göre \Log{G1c}'de
  atomik aksiyomlardan $!B \Sequent !B$ ile $!C \Sequent !C$'nin
  !!{proof}s{}ı~$\pi_1$ ve~$\pi_2$ vardır. Bunları şöyle $!B \land !C
  \Sequent !B \land !C$'nin bir !!a{proof}{}a genişletebiliriz:
  \[
  \AxiomC{}
  \RightLabel{$\pi_1$}
  \Deduce$!B, \fCenter !B$
  \RightLabel{\LeftR{\Weakening}}
  \UnaryInf$!B, !C, \fCenter !B$
  \RightLabel{\LeftR{\land}}
  \UnaryInf$!B \land !C, !C\fCenter !B$
  \RightLabel{\LeftR{\land}}
  \UnaryInf$!B \land !C, !B \land !C \fCenter !B$
  \AxiomC{}
  \RightLabel{$\pi_2$}
  \Deduce$!C, \fCenter !C$
  \RightLabel{\LeftR{\Weakening}}
  \UnaryInf$!B, !C, \fCenter !C$
  \RightLabel{\LeftR{\land}}
  \UnaryInf$!B \land !C, !C \fCenter !C$
  \RightLabel{\LeftR{\land}}
  \UnaryInf$!B \land !C, !B \land !C \fCenter !C$
  \RightLabel{\RightR{\land}}
  \BinaryInf$!B \land !C, !B \land !C \fCenter !B \land !C$
  \RightLabel{\LeftR{\Contraction}}
  \UnaryInf$!B \land !C \fCenter !B \land !C$
  \DisplayProof
  \]

  Şimdi $!A \ident \lforall[x][B(x)]$ olsun. $\lforall[x][B(x)]$ içinde
  geçmeyen bir !!a{constant}~$c$ seçin. O zaman $\depth{B(c)} =
  \depth{!B(x)} < \depth{!A}$ ve tümevarım varsayımına göre \Log{G1c}'de
  atomik aksiyomlardan $!B(c) \Sequent !B(c)$'nin bir !!a{proof}{}ı~$\pi_1$
  vardır. Bunu şöyle $\lforall[x][B(x)] \Sequent \lforall[x][B(x)]$'nin
  bir !!a{proof}{}a genişletebiliriz:
  \[
  \AxiomC{}
  \RightLabel{$\pi_1$}
  \Deduce$!B(c) \fCenter !B(c)$
  \RightLabel{\LeftR{\lforall}}
  \UnaryInf$\lforall[x][!B(x)] \fCenter !B(c)$
  \RightLabel{\RightR{\lforall}}
  \UnaryInf$\lforall[x][!B(x)] \fCenter \lforall[x][!B(x)]$
  \DisplayProof
  \]

  Geri kalan durumlar alıştırma olarak bırakılmıştır.
\end{proof}

\begin{prob}
  $!A \ident (!B \lor !C)$, $!A \ident (!B \lif !C)$ ve $!A \ident
  \lexists[x][B(x)]$ durumlarını ele alarak \olref{prop:G1c-axioms}
  ispatını tamamlayın.
\end{prob}

Amaçlarımız bakımından $!A \ident \lnot !B$ durumunda şu !!{proof}{}ı da
kullanabilirdik:
\[
\AxiomC{}
\RightLabel{$\pi_1$}
\Deduce$!B \fCenter !B$
\RightLabel{\RightR{\lnot}}
\UnaryInf$\fCenter !B, \lnot !B$
\RightLabel{\LeftR{\lnot}}
\UnaryInf$\lnot !B \fCenter \lnot !B$
\DisplayProof
\]
Ancak bunun işlemediği sekant sistemleri vardır. Sezgici sekant hesabı
\Log{G1i}, her sekantın sağ yanında en fazla bir !!{formula} bulunmasını
gerektirmesi dışında \Log{G1c} gibidir. $\Delta = \emptyset$ ise ilk
!!{proof} \Log{G1i}'de de bir !!a{proof} olur; ikincisi olmaz, çünkü bir
sekant sağda hem $!B$ hem de $\lnot !B$'yi içerir.

\Log{G1c}'de bile $!A \ident \lforall[x][!B(x)]$ durumunda kuralların
başka bir sırasını kullanamazdık. Aşağıdaki bir !!a{proof} \emph{değildir}:
\[
\AxiomC{}
\RightLabel{$\pi_1$}
\Deduce$!B(c) \fCenter !B(c)$
\RightLabel{\RightR{\lforall}}
\UnaryInf$!B(c) \fCenter \lforall[x][!B(x)]$
\RightLabel{\LeftR{\lforall}}
\UnaryInf$\lforall[x][!B(x)] \fCenter \lforall[x][!B(x)]$
\DisplayProof
\]
Çünkü $\RightR{\lforall}$ üzerindeki özdeğişken koşulu ihlal edilir.

\begin{prop}\ollabel{prop:G3c-axioms}
$!A$'nın atomik olması \emph{gerekmeksizin}, $!A, \Gamma \Sequent \Delta,
!A$ biçimindeki her sekant \Log{G3c}'de ispatlanabilir.
\end{prop}

\begin{proof}
İspat \olref{prop:G1c-axioms} ispatına çok benzer; ancak tümevarım
varsayımına dikkat etmeliyiz. $n = \depth{!A} = 0$ durumu yine kolaydır.
$n=0$ ise $!A$ atomiktir ve $!A, \Gamma \Sequent \Delta, !A$ bir
\Log{G3c} aksiyomudur.

Şimdi $n>0$ olduğunu varsayalım. Yine durumları ayırırız. Tümevarım
varsayımı şudur: $\depth{!D} < n$ olan bütün $!D$ ile bütün $\Pi$ ve
$\Lambda$ için $\Log{G3c} \Proves !D, \Pi \Sequent \Lambda, !D$.
$!A \ident (!B \land !C)$ durumunu ele alalım. Tümevarım varsayımını iki
kez kullanırız: $!D = !B$, $\Pi = !C, \Gamma$, $\Lambda = \Delta$ alarak
$!B, !C, \Gamma \Sequent \Delta, !B$'nin bir !!a{proof}{}ı~$\pi_1$;
$!D = !C$, $\Pi = !B, \Gamma$, $\Lambda = \Delta$ alarak $!B, !C,
\Gamma \Sequent \Delta, !C$'nin bir !!a{proof}{}ı~$\pi_2$ elde edilir.
$!B \land !C, \Gamma \Sequent \Delta, !B \land !C$'nin !!{proof}{}ı
şöyledir:
\[
\AxiomC{}
\RightLabel{$\pi_1$}
\Deduce$!B, !C, \Gamma \fCenter \Delta, !B$
\RightLabel{\LeftR{\land}}
\UnaryInf$!B \land !C, \Gamma \fCenter \Delta, !B$
\AxiomC{}
\RightLabel{$\pi_2$}
\Deduce$!B, !C, \Gamma \fCenter \Delta, !C$
\RightLabel{\LeftR{\land}}
\UnaryInf$!B \land !C, \Gamma \fCenter \Delta, !C$
\RightLabel{\RightR{\land}}
\BinaryInf$!B \land !C, \Gamma \fCenter \Delta, !B \land !C$
\DisplayProof
\]

$!A \ident \lforall[x][!B(x)]$ için $\lforall[x][!B(x)]$, $\Gamma$ ya da
$\Delta$ içinde geçmeyen bir !!a{constant}~$c$ seçin. $!D = !B(c)$,
$\Pi = \lforall[x][!B(x)], \Gamma$ ve $\Lambda = \Delta$ için tümevarım
varsayımını uygulayın. Böylece $!B(c), \lforall[x][!B(x)], \Gamma
\Sequent \Delta, !B(c)$'nin bir !!a{proof}{}ı~$\pi_1$ elde edilir ve bundan
şunu kurarız:
\[
\AxiomC{}
\RightLabel{$\pi_1$}
\Deduce$!B(c), \lforall[x][!B(x)], \Gamma \fCenter \Delta, !B(c)$
\RightLabel{\LeftR{\lforall}}
\UnaryInf$\lforall[x][!B(x)], \Gamma \fCenter \Delta, !B(c)$
\RightLabel{\RightR{\lforall}}
\UnaryInf$\lforall[x][!B(x)], \Gamma \fCenter \Delta, \lforall[x][!B(x)]$
\DisplayProof
\]
$c$'yi seçme biçimimiz nedeniyle $\RightR{\lforall}$ üzerindeki
özdeğişken koşulu sağlanır.
\end{proof}

\begin{prob}
  $!A \ident (!B \lif !C)$ ve $!A \ident \lexists[x][B(x)]$ durumlarını
  ele alarak \olref{prop:G3c-axioms} ispatını sürdürün.
\end{prob}

\end{document}
